\section{大数定律与长期平均}
\label{sec:05-Probability/08-Limit-Theorems/02}

轮盘赌桌上连出了十次红色。旁边的赌客压低声音对你说：``黑色该来了——概率总要回归的。''

你知道他说错了。轮盘没有记忆，下一转红色和黑色的概率与之前的历史毫无关系。但你又隐约知道另一件事是对的：长期去看，红色的出现比例确实会趋近于它的理论概率 \(18/38 \approx 0.474\)。

这两件事怎么同时成立？如果每次试验独立，为什么长期频率还能稳定到某个数？如果它稳定了，为什么不能用来预测下一次？答案藏在大数定律的精确陈述里——它说的是平均如何趋于常数，而不是下一次如何被过去纠正。

设 \(X_1, X_2, \ldots\) 独立同分布，均值 \(\mu = \E[X_1]\) 存在。样本均值为

\[
\bar X_n = \frac{1}{n}\sum_{i=1}^{n} X_i.
\]

我们关心当 \(n\to\infty\) 时 \(\bar X_n\) 的行为。直观：越多的独立样本应该让平均越可靠。大数定律把这种直觉做成了数学定理——并且在两个不同强度上。

\subsection{弱大数定律：每一次偏离都越来越不可能}

弱大数定律（WLLN）断言

\[
\bar X_n \xrightarrow{P} \mu.
\]

也就是对任何 \(\varepsilon > 0\)，当 \(n\) 充分大时，样本均值偏离真均值超过 \(\varepsilon\) 的概率任意小。

如果额外假设方差 \(\sigma^2 < \infty\)，证明几乎不值一提：

\[
\Var(\bar X_n) = \Var\!\left(\frac{1}{n}\sum_i X_i\right) = \frac{1}{n^2}\sum_i \Var(X_i) = \frac{\sigma^2}{n}.
\]

样本均值的方差以 \(1/n\) 的速度收缩。套上 Chebyshev：

\[
P(|\bar X_n - \mu| \ge \varepsilon) \le \frac{\sigma^2}{n\varepsilon^2} \to 0.
\]

三行，完成。但这短短三行包含了三个独立的事实，混在一起容易漏掉其中一个的微妙之处。拆开看：

\begin{enumerate}
\tightlist
\item
  独立同分布让 \(\Var(\sum_i X_i) = n\sigma^2\)——方差可加但不交叉。若存在正相关，方差会更大；负相关则更小。独立同分布是方差精确为 \(n\sigma^2\) 的保证。
\item
  Chebyshev 把方差的尺度翻译成概率上界。\(\Var(\bar X_n) = \sigma^2/n\) 意味着典型偏差在 \(\sigma/\sqrt{n}\) 尺度上，但 Chebyshev 能证明的只是概率上界为 \(\sigma^2/(n\varepsilon^2)\)——一个 \(1/n\) 而不是 \(1/\sqrt{n}\) 的衰减。尾概率的上界收敛速度比典型偏差的尺度快一档，因为 Chebyshev 用的是二阶矩，天然给出 \(1/\varepsilon^2\) 的因子。
\item
  \(\sigma^2/(n\varepsilon^2) \to 0\) 对任意固定的 \(\varepsilon>0\) 成立，但当 \(\varepsilon\) 随 \(n\) 变化时——比如 \(\varepsilon = 1/\sqrt[3]{n}\)——上界不一定趋于零。这一点在应用 Chebyshev 做一致性证明时常常被忽略。
\end{enumerate}

但有限方差只是这个证明的充分条件，不是弱大数定律本身的必要条件。独立同分布且 \(\E|X_1| < \infty\) 已经足够让同样的依概率收敛结论成立——Khintchine 弱大数定律正是此意。证明需要截断的额外手续：把 \(X_i\) 截到 \([-n, n]\)，在截断变量上使用 Chebyshev（因为截断后方差有限），再证明截断误差对样本平均的贡献依概率趋于零。最后的论证归结为：若 \(\E|X_1|<\infty\)，则 \(\frac{1}{n}\sum_i |X_i|\mathbf{1}_{\{|X_i|>n\}} \xrightarrow{P} 0\)——这一条需要一阶矩的存在性。方差存在的意义是给了你收敛速度：概率上界是 \(O(1/n)\)。如果只是要知道``会不会趋近''而不管速度，一阶矩就够了。

拿一个具体的骰子问题验证一下：掷一枚公平六面骰子，每次结果的均值 \(\mu=3.5\)，方差 \(\sigma^2=35/12\approx 2.917\)。掷 \(n=1000\) 次后，Chebyshev 告诉你：
\[
P(|\bar X_{1000} - 3.5| \ge 0.1) \le \frac{2.917}{1000 \times 0.01} \approx 0.29.
\]
掷 \(n=10000\) 次：
\[
P(|\bar X_{10000} - 3.5| \ge 0.1) \le \frac{2.917}{10000 \times 0.01} \approx 0.029.
\]
每个数量级的样本量把上界压降一个数量级。但这仍是上界——真实的概率通常远小于 Chebyshev 给出的数字，因为骰子结果远比分布无关的 Chebyshev 界所设想的极端情况温和得多。

\subsection{强大数定律：几乎每一条路径都最终忠诚}

强大数定律（SLLN）比弱版本更强：

\[
\bar X_n \xrightarrow{\text{a.s.}} \mu.
\]

除去一个总概率为零的糟糕路径集合之后，你取任意一条具体的无限长样本序列，沿着它一直算运行平均，这个平均值将真正趋近于 \(\mu\)。

强版本和弱版本的差异不在于``收敛快慢''，而在于量词结构。考虑一枚公平硬币（\(p=0.5\)）的无限次抛掷。

\begin{itemize}
\tightlist
\item
  弱定律说：第 \(n\) 次抛完后检查偏差概率——抛 \(10000\) 次，偏离 \(0.01\) 以上的概率已经微乎其微。但不能排除某些具体序列在第 \(10001\) 次又跳远了。弱定律允许坏事件集合随 \(n\) 移动——那些在第 \(n\) 步偏离远的路径，在第 \(n+1\) 步可能回归，也可能换一批新的路径偏出去。
\item
  强定律说：在全空间上抽取一条无限序列，它的运行平均必然收敛到 \(0.5\)，除了一个概率零的例外集（例如永远出正面的序列，其概率为 \(\lim_n (0.5)^n = 0\)，确实属于零概集）。强定律的例外集是全局的——它固定了那批倒霉路径，对于其余所有路径，数列 \(\bar X_1(\omega), \bar X_2(\omega), \ldots\) 就是一个普通的收敛数列。
\end{itemize}

强定律的证明比弱定律复杂得多。关键障碍：单次 Chebyshev 只能压住每一步的边际偏差概率，不足以联合地制服整条路径——因为无穷多个``每一步都大概率接近''的事件同时发生，不是``每一步各自大概率接近''的简单合取。Borel--Cantelli 引理是突破点：它能判定一组事件中``无穷多个发生''的概率是零还是一，从而把逐点概率控制升级为路径收敛。

标准 SLLN 证明路线涉及以下关键步骤：

\begin{enumerate}
\tightlist
\item
  \textbf{截断}。定义 \(Y_k = X_k \cdot \mathbf{1}_{\{|X_k|\le k\}}\)，把每个变量砍到 \([-k,k]\) 内。截断后每个 \(Y_k\) 有界，高阶矩存在。需要论证：砍掉的部分对样本平均的贡献趋于零 a.s.——这需要 \(\E|X_1|<\infty\)（否则截断尾巴的和发散）。
\item
  \textbf{控制截断和的四阶矩}。对截断变量求 \(\E[(Y_k-\E Y_k)^4]\)，利用独立性让交叉项消失，得到上界 \(\le k^2\E[Y_k^2]\) 型的不等式。除以 \(n^4\) 后累加收敛——这就是 Borel--Cantelli 可以处理的 \(1/n^2\) 级衰减。
\item
  \textbf{Borel--Cantelli 到 a.s. 收敛}。因为 \(\sum_n P(|\bar Y_n - \E\bar Y_n| > \varepsilon) < \infty\)（由 Chebyshev + 四阶矩上界保证），第一 Borel--Cantelli 引理推出：对每个 \(\varepsilon>0\)，\(|\bar Y_n - \E\bar Y_n| > \varepsilon\) 只发生有限多次 a.s.。这就得到 \(\bar Y_n - \E\bar Y_n \to 0\) a.s.。
\item
  \textbf{回推原变量}。\(\E Y_k \to \mu\)（截断的期望趋于原期望，由 \(\E|X_1|<\infty\) 保证），且截断误差趋于零 a.s.。组合得到 \(\bar X_n \to \mu\) a.s.。
\end{enumerate}

Kolmogorov 的经典结果用了独立同分布且 \(\E|X_1|<\infty\) 的同一个一阶矩条件，同时给出依概率和几乎处处收敛——这意味着强定律的门槛并没有比弱定律高，只是证明长了数倍。对于非数学专业的读者，记住这个事实即可：在独立同分布下，只要一阶矩存在，样本平均以强且弱两种意义收敛到总体均值。

\subsection{频率解释：概率与长期比例之间的桥}

把 \(X_i\) 设为事件 \(A\) 在第 \(i\) 次独立重复试验中的指示变量：

\[
P(X_i = 1) = p,\qquad P(X_i = 0) = 1-p.
\]

则 \(\bar X_n\) 恰好是前 \(n\) 次试验中 \(A\) 发生的相对频率：

\[
\bar X_n = \frac{\#\{\text{前 }n\text{ 次中 }A\text{ 发生}\}}{n}.
\]

大数定律给出 \(\bar X_n \to p\)（依概率或几乎处处，取决于用的是弱版本还是强版本）。于是概率模型中的数值 \(p\) 与可观察的长期相对频率衔接了起来：概率不是某种神秘的内在倾向，而是``在大量独立重复下频率会稳定到的那个数''——条件是你有一个可重复的、独立同分布的生成机制。

这就是频率学派对概率的诠释核心。它不宣称任何单次试验具有``内在概率''，而是断言：如果你能无限次地重复同一个随机实验，频率极限存在且等于 \(p\)。

但这恰恰是回到赌客谬误的关键。大数定律的运作机制不是``纠偏''，而是``稀释''。

假设你抛了一枚公平硬币 \(10\) 次，得到 \(9\) 次正面、\(1\) 次反面。频率是 \(0.9\)，离 \(0.5\) 差 \(0.4\)。你再抛 \(990\) 次。如果这 \(990\) 次中正面和反面大致各半（正如独立性和 \(p=0.5\) 所预测），加上前 \(10\) 次后总共是 \(9+495=504\) 次正面和 \(1+495=496\) 次反面，频率变成 \(504/1000=0.504\)。接近 \(0.5\) 了。

不是后来的抛掷补偿了前期的偏差——后来的抛掷仍然是公平的，正面概率始终是 \(0.5\)。前期那 \(0.4\) 的绝对偏差（\(9-1=8\) 个正面的净超出）也并没有消失。它只是被后来新增的 \(990\) 次试验产生的庞大分母冲淡了。更精确地，偏差 \(\bar X_n - p\) 的绝对值在以 \(1/\sqrt{n}\) 的量级缩小——不是因为偏差被修正了（净超出 8 始终是 8），而是因为分母 \(n\) 的增长压住了分子，如同 \(8/10 \to 8/1000\)。

这件稀释机制还有一个有趣的推论：回归均值不需要任何``纠偏力''。如果一次考试你考得特别好（远高于你的真实水平），下一次你大概率会回落——不是因为考试变难了，而是因为那次极端高分的一部分来自运气（独立随机误差），而运气没有记忆。大数定律在个体层面也隐约运作：多次观测的平均会消去独立误差项。

\subsection{需要多少样本？一个实用计算}

大数定律只保证最终收敛，不担保任何有限 \(n\) 下的精密度。但在方差有限的条件下，Chebyshev 可以给出一个保守的样本量估计。

回到骰子问题。你想让样本均值与 \(3.5\) 的差距以 \(95\%\) 的概率不超过 \(0.05\)：

\[
P(|\bar X_n - 3.5| \ge 0.05) \le \frac{2.917}{n \cdot 0.0025} \le 0.05.
\]

解得 \(n \ge \frac{2.917}{0.0025\cdot 0.05} = 23336\)。Chebyshev 告诉你需要超过两万次投掷——这显然太保守了，因为 Chebyshev 考虑了最坏的分布形态。实际上，骰子分布相当集中（在 \([1,6]\) 上有界），Hoeffding 不等式给出的样本量远小于此：\(n \ge \frac{\log(2/0.05)}{2\cdot 0.05^2} \approx 740\)。即使是 \(740\) 次投掷，真实覆盖率也远高于 \(95\%\)——CLT 近似会给出约 \(n \ge (1.96\cdot 1.708/0.05)^2 \approx 4500\)（用了标准差），但这是基于正态的渐近值。

这个计算案例揭示了一个核心张力：大数定律不保证快慢，Chebyshev 过于保守，Hoeffding 需要独立有界，CLT 只给渐近。在不同的分析目标下（理论证明 vs 工程估算），你会选择不同的工具——但没有一个单独的结果能同时满足所有需求。

\subsection{大数定律不提供误差形状}

弱大数定律告诉你 \(\bar X_n - \mu \to 0\)。它没有告诉你这个差异的典型大小、分布形状或置信区间。

如果你还知道方差 \(\sigma^2 < \infty\)，你可以多说出一点：\(\bar X_n\) 的标准差是 \(\sigma/\sqrt n\)。这说明典型误差以 \(1/\sqrt n\) 的量级衰减，而不是 \(1/n\)。但就连这条信息都没有从大数定律本身得出——方差是单独假设的，收敛速度需要 Chebyshev（或更精细的工具）。

为了获得误差的分布形状——什么时候能用 Gaussian 近似，什么时候需要更厚的尾部——你需要下一节中心极限定理。

还有一个容易被忽视的区分：一致性（大数定律保证的性质）与无偏性。一个估计量可以在每个有限 \(n\) 都有偏（\(\E[\hat\theta_n] \neq \theta\)），但只要偏差随 \(n\) 趋于零且方差也趋于零，它仍然是一致估计量。例如用 \(\frac{1}{n}\sum (X_i-\bar X_n)^2\) 估计方差——它有偏（期望是 \(\frac{n-1}{n}\sigma^2\)），但偏差 \(-\sigma^2/n \to 0\)，它是一致的。反过来，一个每个 \(n\) 都无偏的估计量，如果方差不随 \(n\) 消失，就不一致。大样本下，偏可以被淹没，波动不行。

\subsection{条件崩塌清单：大数定律不万灵}

大数定律的结论精致地依赖于一组前提。只要一条崩塌，``平均终将稳定''就失效。

\textbf{重尾与无均值}。Cauchy 分布：密度 \(\frac{1}{\pi(1+x^2)}\)，尾部衰减得太慢（对称双尾以 \(1/|x|\) 的速度坠落），导致 \(\E|X| = \infty\)。独立同分布的 Cauchy 变量的样本平均不会趋向任何有限常数——事实上，\(\bar X_n\) 仍然服从 Cauchy 分布，不管 \(n\) 取多大。

这令人震惊：取 \(10\) 个样本的平均和取 \(1\) 个样本，分布的形状一模一样，没有变窄。你可以模拟一万次取一千个 Cauchy 样本的平均，这些平均值不会簇拥在任何有限数附近——偶尔一个极端样本就能把整批平均拉飞。为什么 Cauchy 这么特殊？因为其特征函数是 \(e^{-|t|}\)（没有二阶 Taylor 展开所需的二次项），标准化和的特征函数是 \((e^{-|t|/\sqrt{n}})^n = e^{-|t|}\)，恰好不变。大数定律的前提``\(\E|X_1|<\infty\)''不是一个技术细节，它是结论能够成立的基石。在实际数据中，如果你怀疑尾部很厚（如金融收益的日对数变化、地震震级、互联网流量爆发），在做平均推断之前要高度警觉：均值这个量本身可能不存在，用样本平均去估计它是自欺。此时需要改用中位数或其他对重尾稳健的汇总统计量。

\textbf{强依赖}。若 \(X_i\) 存在长期正相关——比如金融收益率的波动聚类——样本均值的方差不会按 \(1/n\) 缩小。某些平稳过程的样本均值方差甚至可以稳定在非零常数，大数定律完全失效。此时需要遍历定理：在平稳和遍历条件下，时间平均仍收敛到空间平均，但收敛速率不再是 \(1/\sqrt{n}\)，分布的独立性被更弱的混合条件取代。

\textbf{分布漂移}。如果数据生成机制随时间缓慢变化，你所谓的``真正均值''本身就是一个移动靶。样本平均会追踪这个移动目标，但不会收敛到任何一个固定的 \(\mu\)。在线学习和变化点检测正是应对这类情形。

\textbf{选择偏差}。如果你的样本不是从总体中均匀随机抽取的——比如在线调查中只有强烈意见的人才会回应——那么大样本下的平均会稳定地趋向一个错误的数值，而且精度随着 \(n\) 增大而变高，错误也变得越顽固。大数定律让平均值收敛到抽样分布的均值——如果抽样分布本身不是目标总体，那么收敛到的就是一个系统性的错误目标。

\textbf{系统误差}。任何固定的测量偏差——比如秤永远多报 2 克——不会被重复测量消除。\(\bar X_n \to \mu + 2\)，而不是 \(\mu\)。大样本消除随机误差，但对系统误差无动于衷。在物理实验中，随机误差用多次测量来压制，系统误差需要用不同仪器、不同方法交叉验证来暴露和校准。

\textbf{函数的大数定律}。一个容易被忽略的陷阱：\(X_n \to \mu\) a.s. 不直接推出 \(g(X_n) \to g(\mu)\)。需要 \(g\) 在 \(\mu\) 处连续。如果 \(g\) 在 \(\mu\) 处间断——比如 \(g(x)=\mathbf{1}_{x>0}\) 而真实 \(\mu=0\)，微小的随机波动跨过零就会让 \(g(X_n)\) 在 \(0\) 和 \(1\) 间弹跳，平均不收敛到 \(g(0)\)。连续映射定理（CMT）让 \(g\) 连续时依概率/a.s. 收敛保持，但连续性是必要条件。类似地，\(\bar X_n \to \mu\) 不保证 \(1/\bar X_n \to 1/\mu\)（若 \(\mu \neq 0\) 则 \(1/x\) 在 \(\mu\) 处连续，CMT 适用；但若 \(\mu=0\) 则 \(1/x\) 在极限点附近无界，结论崩塌）。

真实数据分析中，核实抽样机制、检查分布稳定性和评估依赖结构，往往比单纯追求更大的样本量更有价值。大数定律替你消灭了随机波动，却不会替你修正设计缺陷。它是一架会自动运行的机器——前提是你要先把它架在正确的地基上。
